% --- Archon physics formalization source begin ---
% archon:physics
% archon:covers PhyXMiniProblems/problem_phyx_mini_0488.lean
% archon:source-report reports/phyx_mini/problem_phyx_mini_0488.source.json

\chapter{Physics problem phyx\_mini\_0488}
\label{ch:PhyXMiniProblems_problem_phyx_mini_0488}

\paragraph{Problem source.}
\#\# Physical scenario

A copper rod and an aluminum rod of the same length and cross-sectional area are attached end to end (figure). The copper end is placed in a furnace maintained at a constant temperature of \textbackslash{}(205\textasciicircum{}\textbackslash{}circ\textbackslash{}mathrm\{C\}\textbackslash{}). The aluminum end is placed in an ice bath held at a constant temperature of \textbackslash{}(0.0\textasciicircum{}\textbackslash{}circ\textbackslash{}mathrm\{C\}\textbackslash{}).

\#\# Question

Calculate the temperature at the point where the two rods are joined.

\#\# Answer choices

- A: A: 100°C

- B: B: 110°C

- C: C: 120°C

- D: D: 130°C

\#\# Figure caption (auxiliary; use the image as primary evidence)

The image shows a horizontal bar composed of two segments: one labeled "Cu" (copper) with a golden hue, and the other labeled "Al" (aluminum) with a silver hue. The copper segment on the left is marked with a temperature of 205°C. The aluminum segment on the right is at 0.0°C. Between them is the notation "T = ?", indicating an unknown temperature at the junction.

\#\# Dataset metadata

Temperature and Heat Transfer; Physical Model Grounding Reasoning, Multi-Formula Reasoning

\paragraph{Recorded answer/context.}
D: D: 130°C

\paragraph{Figure/image path.}
/root/proposal\_for\_physic/science-mango/phyx\_mini\_run/phyx\_data/test\_image/488.png

\paragraph{Formalization target.}
create a compiling Lean file with sorry bodies at `PhyXMiniProblems/problem\_phyx\_mini\_0488.lean`. The Lean declarations must preserve the physical quantities, dimensions or dimensional roles, figure labels, governing-law hypotheses, and final relation expressed by this problem.
Use Mathlib/Physlib names found through LeanExplore where available. If a physics API is missing, introduce faithful local abstractions rather than scalar placeholder aliases.

\begin{theorem}[Physics formalization target]
\label{thm:physics:phyx_mini_0488:target}
\lean{PhyXMiniProblems.ProblemPhyXMini0488.problem_phyx_mini_0488}
\uses{def:physics:phyx-mini-0488:phyxminiproblems-problemphyxmini0488-compositerodsetup, def:physics:phyx-mini-0488:phyxminiproblems-problemphyxmini0488-matchescompositerodscenario, def:physics:phyx-mini-0488:phyxminiproblems-problemphyxmini0488-matchessuppliedfigure, def:physics:phyx-mini-0488:phyxminiproblems-problemphyxmini0488-matchesproblemreadouts, def:physics:phyx-mini-0488:phyxminiproblems-problemphyxmini0488-usestextbookthermalconductivities, def:physics:phyx-mini-0488:phyxminiproblems-problemphyxmini0488-hasphysicalcompositerodparameters, def:physics:phyx-mini-0488:phyxminiproblems-problemphyxmini0488-satisfiesperfectthermalcontactlaw, def:physics:phyx-mini-0488:phyxminiproblems-problemphyxmini0488-satisfiessteadyseriesfourierconduction, def:physics:phyx-mini-0488:phyxminiproblems-problemphyxmini0488-junctiontemperatureindegreescelsius, def:physics:phyx-mini-0488:phyxminiproblems-problemphyxmini0488-matchesanswerchoice}
The assigned autoformalize agent should translate this physics problem into Lean declarations in the covered file, with theorem and lemma proof bodies written as `by sorry`.
\end{theorem}
\begin{proof}
This is an autoformalization task, not a proof task. Produce statements that can later be proved by the physics prover without weakening the physical model.
\end{proof}

% --- Archon declaration topology begin ---
\section{Declaration topology}
\begin{definition}[Length Quantity]
  \label{def:physics:phyx-mini-0488:phyxminiproblems-problemphyxmini0488-lengthquantity}
  \lean{PhyXMiniProblems.ProblemPhyXMini0488.LengthQuantity}
  A nonnegative physical length, independent of a choice of unit
\end{definition}
\begin{proof}
  This is the declared model definition of Length Quantity.
\end{proof}
\begin{definition}[Area Quantity]
  \label{def:physics:phyx-mini-0488:phyxminiproblems-problemphyxmini0488-areaquantity}
  \lean{PhyXMiniProblems.ProblemPhyXMini0488.AreaQuantity}
  A nonnegative physical cross-sectional area
\end{definition}
\begin{proof}
  This is the declared model definition of Area Quantity.
\end{proof}
\begin{definition}[Heat Current Quantity]
  \label{def:physics:phyx-mini-0488:phyxminiproblems-problemphyxmini0488-heatcurrentquantity}
  \lean{PhyXMiniProblems.ProblemPhyXMini0488.HeatCurrentQuantity}
  A nonnegative heat-transfer rate, with the physical dimension of power
\end{definition}
\begin{proof}
  This is the declared model definition of Heat Current Quantity.
\end{proof}
\begin{definition}[Thermal Conductivity Quantity]
  \label{def:physics:phyx-mini-0488:phyxminiproblems-problemphyxmini0488-thermalconductivityquantity}
  \lean{PhyXMiniProblems.ProblemPhyXMini0488.ThermalConductivityQuantity}
  Thermal conductivity, whose coherent SI unit is watt per metre-kelvin
\end{definition}
\begin{proof}
  This is the declared model definition of Thermal Conductivity Quantity.
\end{proof}
\begin{definition}[length Readout]
  \label{def:physics:phyx-mini-0488:phyxminiproblems-problemphyxmini0488-lengthreadout}
  \lean{PhyXMiniProblems.ProblemPhyXMini0488.lengthReadout}
  \uses{def:physics:phyx-mini-0488:phyxminiproblems-problemphyxmini0488-lengthquantity}
  Read a physical length in a selected length unit
\end{definition}
\begin{proof}
  This is the declared model definition of length Readout.
\end{proof}
\begin{definition}[area Readout]
  \label{def:physics:phyx-mini-0488:phyxminiproblems-problemphyxmini0488-areareadout}
  \lean{PhyXMiniProblems.ProblemPhyXMini0488.areaReadout}
  \uses{def:physics:phyx-mini-0488:phyxminiproblems-problemphyxmini0488-areaquantity}
  Read an area in the square of a selected length unit
\end{definition}
\begin{proof}
  This is the declared model definition of area Readout.
\end{proof}
\begin{definition}[heat Current Readout]
  \label{def:physics:phyx-mini-0488:phyxminiproblems-problemphyxmini0488-heatcurrentreadout}
  \lean{PhyXMiniProblems.ProblemPhyXMini0488.heatCurrentReadout}
  \uses{def:physics:phyx-mini-0488:phyxminiproblems-problemphyxmini0488-heatcurrentquantity}
  Read a heat current in the coherent unit determined by the base units
\end{definition}
\begin{proof}
  This is the declared model definition of heat Current Readout.
\end{proof}
\begin{definition}[thermal Conductivity Readout]
  \label{def:physics:phyx-mini-0488:phyxminiproblems-problemphyxmini0488-thermalconductivityreadout}
  \lean{PhyXMiniProblems.ProblemPhyXMini0488.thermalConductivityReadout}
  \uses{def:physics:phyx-mini-0488:phyxminiproblems-problemphyxmini0488-thermalconductivityquantity}
  Read thermal conductivity in the coherent unit determined by the base units
\end{definition}
\begin{proof}
  This is the declared model definition of thermal Conductivity Readout.
\end{proof}
\begin{definition}[length In Meters]
  \label{def:physics:phyx-mini-0488:phyxminiproblems-problemphyxmini0488-lengthinmeters}
  \lean{PhyXMiniProblems.ProblemPhyXMini0488.lengthInMeters}
  \uses{def:physics:phyx-mini-0488:phyxminiproblems-problemphyxmini0488-lengthquantity, def:physics:phyx-mini-0488:phyxminiproblems-problemphyxmini0488-lengthreadout}
  Metre readout of a physical rod length
\end{definition}
\begin{proof}
  This is the declared model definition of length In Meters.
\end{proof}
\begin{definition}[area In Square Meters]
  \label{def:physics:phyx-mini-0488:phyxminiproblems-problemphyxmini0488-areainsquaremeters}
  \lean{PhyXMiniProblems.ProblemPhyXMini0488.areaInSquareMeters}
  \uses{def:physics:phyx-mini-0488:phyxminiproblems-problemphyxmini0488-areaquantity, def:physics:phyx-mini-0488:phyxminiproblems-problemphyxmini0488-areareadout}
  Square-metre readout of a physical cross-sectional area
\end{definition}
\begin{proof}
  This is the declared model definition of area In Square Meters.
\end{proof}
\begin{definition}[heat Current In Watts]
  \label{def:physics:phyx-mini-0488:phyxminiproblems-problemphyxmini0488-heatcurrentinwatts}
  \lean{PhyXMiniProblems.ProblemPhyXMini0488.heatCurrentInWatts}
  \uses{def:physics:phyx-mini-0488:phyxminiproblems-problemphyxmini0488-heatcurrentquantity, def:physics:phyx-mini-0488:phyxminiproblems-problemphyxmini0488-heatcurrentreadout}
  Watt readout of a physical axial heat current
\end{definition}
\begin{proof}
  This is the declared model definition of heat Current In Watts.
\end{proof}
\begin{definition}[thermal Conductivity In Watts Per Meter Kelvin]
  \label{def:physics:phyx-mini-0488:phyxminiproblems-problemphyxmini0488-thermalconductivityinwattspermeterkelvin}
  \lean{PhyXMiniProblems.ProblemPhyXMini0488.thermalConductivityInWattsPerMeterKelvin}
  \uses{def:physics:phyx-mini-0488:phyxminiproblems-problemphyxmini0488-thermalconductivityquantity, def:physics:phyx-mini-0488:phyxminiproblems-problemphyxmini0488-thermalconductivityreadout}
  Watt-per-metre-kelvin readout of a physical conductivity
\end{definition}
\begin{proof}
  This is the declared model definition of thermal Conductivity In Watts Per Meter Kelvin.
\end{proof}
\begin{definition}[Celsius Temperature Reading]
  \label{def:physics:phyx-mini-0488:phyxminiproblems-problemphyxmini0488-celsiustemperaturereading}
  \lean{PhyXMiniProblems.ProblemPhyXMini0488.CelsiusTemperatureReading}
  Physlib models absolute temperature in a zero-preserving scale, while Celsius is affine. This record retains the absolute physical quantity and explicitly calibrates the scalar thermometer readout to kelvins
\end{definition}
\begin{proof}
  This is the declared model definition of Celsius Temperature Reading.
\end{proof}
\begin{definition}[Rod Section]
  \label{def:physics:phyx-mini-0488:phyxminiproblems-problemphyxmini0488-rodsection}
  \lean{PhyXMiniProblems.ProblemPhyXMini0488.RodSection}
  The two end-to-end sections shown in the supplied image
\end{definition}
\begin{proof}
  This is the declared model definition of Rod Section.
\end{proof}
\begin{definition}[Rod Material]
  \label{def:physics:phyx-mini-0488:phyxminiproblems-problemphyxmini0488-rodmaterial}
  \lean{PhyXMiniProblems.ProblemPhyXMini0488.RodMaterial}
  The two materials named in the problem
\end{definition}
\begin{proof}
  This is the declared model definition of Rod Material.
\end{proof}
\begin{definition}[Thermal Location]
  \label{def:physics:phyx-mini-0488:phyxminiproblems-problemphyxmini0488-thermallocation}
  \lean{PhyXMiniProblems.ProblemPhyXMini0488.ThermalLocation}
  Thermally distinct locations from the furnace to the ice bath
\end{definition}
\begin{proof}
  This is the declared model definition of Thermal Location.
\end{proof}
\begin{definition}[Thermal Contact]
  \label{def:physics:phyx-mini-0488:phyxminiproblems-problemphyxmini0488-thermalcontact}
  \lean{PhyXMiniProblems.ProblemPhyXMini0488.ThermalContact}
  The three thermal interfaces along the series heat-flow path
\end{definition}
\begin{proof}
  This is the declared model definition of Thermal Contact.
\end{proof}
\begin{definition}[Contact Quality]
  \label{def:physics:phyx-mini-0488:phyxminiproblems-problemphyxmini0488-contactquality}
  \lean{PhyXMiniProblems.ProblemPhyXMini0488.ContactQuality}
  Quality of a thermal interface
\end{definition}
\begin{proof}
  This is the declared model definition of Contact Quality.
\end{proof}
\begin{definition}[Rod Arrangement]
  \label{def:physics:phyx-mini-0488:phyxminiproblems-problemphyxmini0488-rodarrangement}
  \lean{PhyXMiniProblems.ProblemPhyXMini0488.RodArrangement}
  Mechanical arrangement of the two rods
\end{definition}
\begin{proof}
  This is the declared model definition of Rod Arrangement.
\end{proof}
\begin{definition}[Thermal Regime]
  \label{def:physics:phyx-mini-0488:phyxminiproblems-problemphyxmini0488-thermalregime}
  \lean{PhyXMiniProblems.ProblemPhyXMini0488.ThermalRegime}
  Time dependence of the temperature profile
\end{definition}
\begin{proof}
  This is the declared model definition of Thermal Regime.
\end{proof}
\begin{definition}[Rod Side Condition]
  \label{def:physics:phyx-mini-0488:phyxminiproblems-problemphyxmini0488-rodsidecondition}
  \lean{PhyXMiniProblems.ProblemPhyXMini0488.RodSideCondition}
  Idealized lateral boundary condition of the rods
\end{definition}
\begin{proof}
  This is the declared model definition of Rod Side Condition.
\end{proof}
\begin{definition}[Reservoir Kind]
  \label{def:physics:phyx-mini-0488:phyxminiproblems-problemphyxmini0488-reservoirkind}
  \lean{PhyXMiniProblems.ProblemPhyXMini0488.ReservoirKind}
  The two constant-temperature devices named in the prose
\end{definition}
\begin{proof}
  This is the declared model definition of Reservoir Kind.
\end{proof}
\begin{definition}[contact Locations]
  \label{def:physics:phyx-mini-0488:phyxminiproblems-problemphyxmini0488-contactlocations}
  \lean{PhyXMiniProblems.ProblemPhyXMini0488.contactLocations}
  \uses{def:physics:phyx-mini-0488:phyxminiproblems-problemphyxmini0488-thermallocation, def:physics:phyx-mini-0488:phyxminiproblems-problemphyxmini0488-thermalcontact}
  The locations joined by each thermal contact
\end{definition}
\begin{proof}
  This is the declared model definition of contact Locations.
\end{proof}
\begin{definition}[section Hot Location]
  \label{def:physics:phyx-mini-0488:phyxminiproblems-problemphyxmini0488-sectionhotlocation}
  \lean{PhyXMiniProblems.ProblemPhyXMini0488.sectionHotLocation}
  \uses{def:physics:phyx-mini-0488:phyxminiproblems-problemphyxmini0488-rodsection, def:physics:phyx-mini-0488:phyxminiproblems-problemphyxmini0488-thermallocation}
  Hot-side location of each rod section
\end{definition}
\begin{proof}
  This is the declared model definition of section Hot Location.
\end{proof}
\begin{definition}[section Cold Location]
  \label{def:physics:phyx-mini-0488:phyxminiproblems-problemphyxmini0488-sectioncoldlocation}
  \lean{PhyXMiniProblems.ProblemPhyXMini0488.sectionColdLocation}
  \uses{def:physics:phyx-mini-0488:phyxminiproblems-problemphyxmini0488-rodsection, def:physics:phyx-mini-0488:phyxminiproblems-problemphyxmini0488-thermallocation}
  Cold-side location of each rod section
\end{definition}
\begin{proof}
  This is the declared model definition of section Cold Location.
\end{proof}
\begin{definition}[Figure Object]
  \label{def:physics:phyx-mini-0488:phyxminiproblems-problemphyxmini0488-figureobject}
  \lean{PhyXMiniProblems.ProblemPhyXMini0488.FigureObject}
  Physical objects visibly represented in the supplied raster
\end{definition}
\begin{proof}
  This is the declared model definition of Figure Object.
\end{proof}
\begin{definition}[Figure Label]
  \label{def:physics:phyx-mini-0488:phyxminiproblems-problemphyxmini0488-figurelabel}
  \lean{PhyXMiniProblems.ProblemPhyXMini0488.FigureLabel}
  Literal labels and numerical marks printed in the supplied raster
\end{definition}
\begin{proof}
  This is the declared model definition of Figure Label.
\end{proof}
\begin{definition}[Composite Rod Figure]
  \label{def:physics:phyx-mini-0488:phyxminiproblems-problemphyxmini0488-compositerodfigure}
  \lean{PhyXMiniProblems.ProblemPhyXMini0488.CompositeRodFigure}
  \uses{def:physics:phyx-mini-0488:phyxminiproblems-problemphyxmini0488-figureobject, def:physics:phyx-mini-0488:phyxminiproblems-problemphyxmini0488-figurelabel}
  Qualitative and label-level evidence transcribed from the primary image
\end{definition}
\begin{proof}
  This is the declared model definition of Composite Rod Figure.
\end{proof}
\begin{definition}[Composite Rod Setup]
  \label{def:physics:phyx-mini-0488:phyxminiproblems-problemphyxmini0488-compositerodsetup}
  \lean{PhyXMiniProblems.ProblemPhyXMini0488.CompositeRodSetup}
  \uses{def:physics:phyx-mini-0488:phyxminiproblems-problemphyxmini0488-lengthquantity, def:physics:phyx-mini-0488:phyxminiproblems-problemphyxmini0488-areaquantity, def:physics:phyx-mini-0488:phyxminiproblems-problemphyxmini0488-heatcurrentquantity, def:physics:phyx-mini-0488:phyxminiproblems-problemphyxmini0488-thermalconductivityquantity, def:physics:phyx-mini-0488:phyxminiproblems-problemphyxmini0488-celsiustemperaturereading, def:physics:phyx-mini-0488:phyxminiproblems-problemphyxmini0488-rodsection, def:physics:phyx-mini-0488:phyxminiproblems-problemphyxmini0488-rodmaterial, def:physics:phyx-mini-0488:phyxminiproblems-problemphyxmini0488-thermallocation, def:physics:phyx-mini-0488:phyxminiproblems-problemphyxmini0488-thermalcontact, def:physics:phyx-mini-0488:phyxminiproblems-problemphyxmini0488-contactquality, def:physics:phyx-mini-0488:phyxminiproblems-problemphyxmini0488-rodarrangement, def:physics:phyx-mini-0488:phyxminiproblems-problemphyxmini0488-thermalregime, def:physics:phyx-mini-0488:phyxminiproblems-problemphyxmini0488-rodsidecondition, def:physics:phyx-mini-0488:phyxminiproblems-problemphyxmini0488-reservoirkind, def:physics:phyx-mini-0488:phyxminiproblems-problemphyxmini0488-compositerodfigure}
  Independent physical data for the apparatus. In particular, neither junction temperature is defined from 130 °C or from any answer-choice value
\end{definition}
\begin{proof}
  This is the declared model definition of Composite Rod Setup.
\end{proof}
\begin{definition}[Matches Composite Rod Scenario]
  \label{def:physics:phyx-mini-0488:phyxminiproblems-problemphyxmini0488-matchescompositerodscenario}
  \lean{PhyXMiniProblems.ProblemPhyXMini0488.MatchesCompositeRodScenario}
  \uses{def:physics:phyx-mini-0488:phyxminiproblems-problemphyxmini0488-compositerodsetup}
  Categorical apparatus and idealized thermal regime used by the model
\end{definition}
\begin{proof}
  This is the declared model definition of Matches Composite Rod Scenario.
\end{proof}
\begin{definition}[Matches Supplied Figure]
  \label{def:physics:phyx-mini-0488:phyxminiproblems-problemphyxmini0488-matchessuppliedfigure}
  \lean{PhyXMiniProblems.ProblemPhyXMini0488.MatchesSuppliedFigure}
  \uses{def:physics:phyx-mini-0488:phyxminiproblems-problemphyxmini0488-compositerodsetup}
  Primary-image evidence. The junction is recorded only as the unknown label T = ?; no numerical junction temperature occurs in this interface
\end{definition}
\begin{proof}
  This is the declared model definition of Matches Supplied Figure.
\end{proof}
\begin{definition}[Matches Problem Readouts]
  \label{def:physics:phyx-mini-0488:phyxminiproblems-problemphyxmini0488-matchesproblemreadouts}
  \lean{PhyXMiniProblems.ProblemPhyXMini0488.MatchesProblemReadouts}
  \uses{def:physics:phyx-mini-0488:phyxminiproblems-problemphyxmini0488-compositerodsetup}
  Geometry and endpoint temperature readouts stated in the problem. The common length and area need no numerical values because they cancel in the series steady-state balance
\end{definition}
\begin{proof}
  This is the declared model definition of Matches Problem Readouts.
\end{proof}
\begin{definition}[Uses Textbook Thermal Conductivities]
  \label{def:physics:phyx-mini-0488:phyxminiproblems-problemphyxmini0488-usestextbookthermalconductivities}
  \lean{PhyXMiniProblems.ProblemPhyXMini0488.UsesTextbookThermalConductivities}
  \uses{def:physics:phyx-mini-0488:phyxminiproblems-problemphyxmini0488-thermalconductivityinwattspermeterkelvin, def:physics:phyx-mini-0488:phyxminiproblems-problemphyxmini0488-compositerodsetup}
  The source does not print the material table required for a numerical answer. The values 385 W/(m K) for copper and 205 W/(m K) for aluminum are exposed as a separate textbook calibration rather than hidden in the target
\end{definition}
\begin{proof}
  This is the declared model definition of Uses Textbook Thermal Conductivities.
\end{proof}
\begin{definition}[Has Physical Composite Rod Parameters]
  \label{def:physics:phyx-mini-0488:phyxminiproblems-problemphyxmini0488-hasphysicalcompositerodparameters}
  \lean{PhyXMiniProblems.ProblemPhyXMini0488.HasPhysicalCompositeRodParameters}
  \uses{def:physics:phyx-mini-0488:phyxminiproblems-problemphyxmini0488-lengthinmeters, def:physics:phyx-mini-0488:phyxminiproblems-problemphyxmini0488-areainsquaremeters, def:physics:phyx-mini-0488:phyxminiproblems-problemphyxmini0488-thermalconductivityinwattspermeterkelvin, def:physics:phyx-mini-0488:phyxminiproblems-problemphyxmini0488-compositerodsetup}
  Positivity and hot-to-cold ordering required by the conduction model
\end{definition}
\begin{proof}
  This is the declared model definition of Has Physical Composite Rod Parameters.
\end{proof}
\begin{definition}[Satisfies Perfect Thermal Contact Law]
  \label{def:physics:phyx-mini-0488:phyxminiproblems-problemphyxmini0488-satisfiesperfectthermalcontactlaw}
  \lean{PhyXMiniProblems.ProblemPhyXMini0488.SatisfiesPerfectThermalContactLaw}
  \uses{def:physics:phyx-mini-0488:phyxminiproblems-problemphyxmini0488-contactlocations, def:physics:phyx-mini-0488:phyxminiproblems-problemphyxmini0488-compositerodsetup}
  A perfect contact equates the physical temperature readings on its two sides. This general interface contains no numerical junction-temperature value
\end{definition}
\begin{proof}
  This is the declared model definition of Satisfies Perfect Thermal Contact Law.
\end{proof}
\begin{definition}[Satisfies Steady Series Fourier Conduction]
  \label{def:physics:phyx-mini-0488:phyxminiproblems-problemphyxmini0488-satisfiessteadyseriesfourierconduction}
  \lean{PhyXMiniProblems.ProblemPhyXMini0488.SatisfiesSteadySeriesFourierConduction}
  \uses{def:physics:phyx-mini-0488:phyxminiproblems-problemphyxmini0488-lengthinmeters, def:physics:phyx-mini-0488:phyxminiproblems-problemphyxmini0488-areainsquaremeters, def:physics:phyx-mini-0488:phyxminiproblems-problemphyxmini0488-heatcurrentinwatts, def:physics:phyx-mini-0488:phyxminiproblems-problemphyxmini0488-thermalconductivityinwattspermeterkelvin, def:physics:phyx-mini-0488:phyxminiproblems-problemphyxmini0488-sectionhotlocation, def:physics:phyx-mini-0488:phyxminiproblems-problemphyxmini0488-sectioncoldlocation, def:physics:phyx-mini-0488:phyxminiproblems-problemphyxmini0488-compositerodsetup}
  Steady one-dimensional Fourier conduction in each homogeneous rod section: heat current = conductivity * area * temperature drop / length. The same independent axial current passes through both series sections, and the ideal side condition makes lateral loss zero. These are governing laws, not a restatement of the requested junction temperature
\end{definition}
\begin{proof}
  This is the declared model definition of Satisfies Steady Series Fourier Conduction.
\end{proof}
\begin{definition}[junction Temperature In Degrees Celsius]
  \label{def:physics:phyx-mini-0488:phyxminiproblems-problemphyxmini0488-junctiontemperatureindegreescelsius}
  \lean{PhyXMiniProblems.ProblemPhyXMini0488.junctionTemperatureInDegreesCelsius}
  \uses{def:physics:phyx-mini-0488:phyxminiproblems-problemphyxmini0488-compositerodsetup}
  Celsius readout on the copper side of the perfect Cu--Al junction
\end{definition}
\begin{proof}
  This is the declared model definition of junction Temperature In Degrees Celsius.
\end{proof}
\begin{theorem}[junction temperature from conductivities]
  \label{thm:physics:phyx-mini-0488:phyxminiproblems-problemphyxmini0488-junction-temperature-from-conductivities}
  \lean{PhyXMiniProblems.ProblemPhyXMini0488.junction_temperature_from_conductivities}
  \uses{def:physics:phyx-mini-0488:phyxminiproblems-problemphyxmini0488-thermalconductivityinwattspermeterkelvin, def:physics:phyx-mini-0488:phyxminiproblems-problemphyxmini0488-compositerodsetup, def:physics:phyx-mini-0488:phyxminiproblems-problemphyxmini0488-matchescompositerodscenario, def:physics:phyx-mini-0488:phyxminiproblems-problemphyxmini0488-matchesproblemreadouts, def:physics:phyx-mini-0488:phyxminiproblems-problemphyxmini0488-hasphysicalcompositerodparameters, def:physics:phyx-mini-0488:phyxminiproblems-problemphyxmini0488-satisfiesperfectthermalcontactlaw, def:physics:phyx-mini-0488:phyxminiproblems-problemphyxmini0488-satisfiessteadyseriesfourierconduction, def:physics:phyx-mini-0488:phyxminiproblems-problemphyxmini0488-junctiontemperatureindegreescelsius}
  The source itself identifies the two materials but does not print a numerical conductivity table. Before adding any textbook calibration, steady series conduction therefore determines the junction temperature symbolically as T = 205 * k\_Cu / (k\_Cu + k\_Al). This is the strongest numerical-readout conclusion supported by the stated apparatus, endpoint temperatures, and governing laws alone. In particular, neither a conductivity value nor an answer-choice value is a premise here
\end{theorem}
\begin{proof}
  The conclusion follows from the governing relations and figure readouts named in the statement of junction temperature from conductivities.
\end{proof}
\begin{definition}[Answer Choice]
  \label{def:physics:phyx-mini-0488:phyxminiproblems-problemphyxmini0488-answerchoice}
  \lean{PhyXMiniProblems.ProblemPhyXMini0488.AnswerChoice}
  Labels of the four answer choices printed in the problem
\end{definition}
\begin{proof}
  This is the declared model definition of Answer Choice.
\end{proof}
\begin{definition}[displayed Junction Temperature Celsius]
  \label{def:physics:phyx-mini-0488:phyxminiproblems-problemphyxmini0488-displayedjunctiontemperaturecelsius}
  \lean{PhyXMiniProblems.ProblemPhyXMini0488.displayedJunctionTemperatureCelsius}
  \uses{def:physics:phyx-mini-0488:phyxminiproblems-problemphyxmini0488-answerchoice}
  Celsius value printed beside each supplied answer label
\end{definition}
\begin{proof}
  This is the declared model definition of displayed Junction Temperature Celsius.
\end{proof}
\begin{definition}[recorded Dataset Answer]
  \label{def:physics:phyx-mini-0488:phyxminiproblems-problemphyxmini0488-recordeddatasetanswer}
  \lean{PhyXMiniProblems.ProblemPhyXMini0488.recordedDatasetAnswer}
  \uses{def:physics:phyx-mini-0488:phyxminiproblems-problemphyxmini0488-answerchoice}
  Dataset metadata recording the supplied answer label. It is not a premise
\end{definition}
\begin{proof}
  This is the declared model definition of recorded Dataset Answer.
\end{proof}
\begin{definition}[Matches Answer Choice]
  \label{def:physics:phyx-mini-0488:phyxminiproblems-problemphyxmini0488-matchesanswerchoice}
  \lean{PhyXMiniProblems.ProblemPhyXMini0488.MatchesAnswerChoice}
  \uses{def:physics:phyx-mini-0488:phyxminiproblems-problemphyxmini0488-compositerodsetup, def:physics:phyx-mini-0488:phyxminiproblems-problemphyxmini0488-junctiontemperatureindegreescelsius, def:physics:phyx-mini-0488:phyxminiproblems-problemphyxmini0488-answerchoice, def:physics:phyx-mini-0488:phyxminiproblems-problemphyxmini0488-displayedjunctiontemperaturecelsius}
  Agreement with a temperature displayed to the nearest ten degrees Celsius
\end{definition}
\begin{proof}
  This is the declared model definition of Matches Answer Choice.
\end{proof}
% --- Archon declaration topology end ---
% --- Archon physics formalization source end ---
